\section{Cardinalities and minimum orders}\label{sec:searches}\label{sec:table}

We count connected graphs up to isomorphism. For orders $3$ through
$13$ the universe is unrestricted; at order $14$ it consists of the
connected $K_4$-free graphs. The exact decisions in
Section~\ref{sec:algorithms}, applied to representatives of these
universes, give the directly computed columns of
Table~\ref{tab:card}: connected graphs, Euler characteristic one,
$\mathbb F_2$-acyclicity, $\IS$, $\ISV$, the four dismantling levels
$0$ through $3$, and cones. The columns for $\I$ and integral
acyclicity are deduced below by classifying the remaining acyclic
graphs.

The finite enumerations and the resulting cardinalities in this
section are computer-assisted. The accompanying material records
the input universes, comparisons and finite verifications, with a map
from the results here to their supporting files.
The mathematical procedures establish what is counted; the retained
execution records supply their numerical results. The finite enumerations
use the generation and isomorphism software of McKay and Piperno
\cite{McKayPiperno2014}; the versions and generation restrictions are
recorded with the data.

\subsection{The first acyclic exceptions}

The directly computed acyclic and strong columns agree through order
$12$. At order $13$ their difference is three. 

\begin{theorem}\label{thm:classification13}
Up to isomorphism, the only $\mathbb F_2$-acyclic graphs on at most
$13$ vertices outside $\IS$ are $\Gt,\Gt',\Gt''$. Every graph in
$\I$ on at most $13$ vertices belongs to $\IS$, and $\Gt$ is the
only stuck graph of order $13$.
\end{theorem}
\begin{proof}
Every $\mathbb F_2$-acyclic graph is connected, so the connected census
contains both classes. Every strong graph is $\mathbb F_2$-acyclic, so equality of the two
finite cardinalities through order $12$ gives equality of the classes
in those orders. The cases of orders one and two are immediate.
At order $13$ the difference is
\[
 \RowThirteenAcyclic-\RowThirteenIS=3.
\]
Theorem~\ref{thm:A} supplies three acyclic non-strong graphs with
different edge counts, hence pairwise non-isomorphic. They exhaust
the difference. Their nonzero integral homology excludes them from
$\I$, which lies between $\IS$ and $\mathcal A_{\mathbb F_2}$.
Finally, the link calculations in Table~\ref{tab:witness-links}
show that $\Gt$ is stuck, whereas $\Gt'$ and $\Gt''$ each admit the
edge deletion reversing the last specified addition. Every stuck graph
is acyclic and non-strong, so there are no others at this order.
\end{proof}

\subsection{The \texorpdfstring{$K_4$}{K4}-free graphs of order fourteen}

The following classification supplies the positive exceptions in
the last row of the table. Its reduction uses only the complete
classification through order $13$ and the local $K_4$-free criterion.

\begin{theorem}\label{thm:classification14}
There are $51$ isomorphism classes of stuck $K_4$-free graphs of
order $14$. Exactly three lie in $\I$, namely $\Ga,\Gb,\Gc$;
each of the other $48$ has first integral homology $\ZZ/3$ and
zero higher homology. There are exactly eight $K_4$-free graphs
of order $14$ in $\I\setminus\IS$, up to isomorphism. Their
minimum number of edges is $44$, attained by $\Ga$.
\end{theorem}
\begin{proof}
First, enumerate the connected $K_4$-free graphs of orders two through
seven and apply Proposition~\ref{rem:K4-local}. No graph is almost
stuck. This check examines, respectively, $1,2,5,17,82,536$
isomorphism classes; it includes graphs with no admissible vertex.
If a stuck $K_4$-free graph of order at most $14$ had a cut vertex,
Lemma~\ref{lem:2conn} would give two almost stuck pieces with
orders summing to at most $15$. One would have order at most seven,
a contradiction. The seeds are therefore $2$-connected and satisfy
$4\le\delta\le\Delta\le12$, by Lemma~\ref{lem:degrees}.

Enumerating the $K_4$-free graphs of order $14$ with these restrictions,
and testing $\mathbb F_2$-homology and the local neighbourhood conditions, gives
$51$ pairwise non-isomorphic stuck graphs. Theorem~\ref{thm:B}
and Proposition~\ref{prop:other-seeds} prove that $\Ga,\Gb,\Gc$
lie in $\I$. Integral homology of the
other $48$, computed by~\eqref{eq:integral-homology}, has free
rank zero in degree one, and $\partial_2$ has a single nonunit Smith
factor, equal to $3$. Thus $H_1\cong\ZZ/3$. Their clique complexes
are connected, have dimension at most two and Euler characteristic
one. Hence $\widetilde H_0=0$, and $H_2$, which is free, has rank
zero. All higher homology therefore vanishes. In particular,
none of these $48$ seeds lies in $\I$.
The list of seeds, together with its homology
and the finite checks, is included in the accompanying repository.

By Theorem~\ref{thm:classification13}, a graph in
$\I\setminus\IS$ has at least $14$ vertices. Apply
Lemma~\ref{lem:stuck} to the $K_4$-free class: every $K_4$-free
member of $\I\setminus\IS$ of order $14$ reduces by edge deletions
to a stuck graph in $\I$.
It therefore arises from one of $\Ga,\Gb,\Gc$. Their pruned
edge-addition closures, restricted to $K_4$-free graphs, have the
following numbers of non-strong classes, grouped by edge count:
\[
\begin{array}{c|rrrr|r}
 \text{seed} & 44 & 45 & 46 & 47 & \text{total}\\ \hline
 \Ga & 1 & 2 & 1 & 0 & 4\\
 \Gb & 0 & 1 & 1 & 1 & 3\\
 \Gc & 0 & 1 & 0 & 0 & 1
\end{array}
\]
The three sets are pairwise disjoint up to isomorphism, so their
union has eight classes. Proposition~\ref{prop:closure-correct}
shows that no non-strong target is lost. Every retained graph belongs
to $\I$, since the additions are I-transformations; every possible
target has now been included. This proves the count of eight.

The three contractible seeds have $44,45,45$ edges. Edge additions
cannot decrease these numbers, so every target has at least $44$
edges. The graph $\Ga$ attains the bound.
\end{proof}

The eight graphs have one-edge extensions in $\IS$, each obtained by
an admissible addition. Thus their membership in $\I$ can also be
checked individually by~\eqref{eq:I-path}; the full lists and
reductions are supplied with the repository.

\subsection{Adjudication of the integral and contractible columns}

The coefficient distinction is essential to the first conjecture.
Within each finite universe, the inclusions
\begin{equation}\label{eq:table-chain}
 \IS\subseteq\I\subseteq\mathcal A_{\ZZ}
       \subseteq\mathcal A_{\mathbb F_2}
\end{equation}
leave only the acyclic non-strong graphs to be decided. Through
order $12$ there are none. At order $13$ all three have integral
torsion, as established above.

At order $14$, let $G$ be $K_4$-free, $\mathbb F_2$-acyclic and
outside $\IS$. An admissible deletion preserves its homotopy type
and cannot leave a strong graph, since that would give a strong
reduction of $G$. Delete admissible edges until either a vertex
deletion is possible or no admissible element remains. In the first
case the vertex deletion gives an acyclic non-strong graph of order
$13$. Theorem~\ref{thm:classification13} identifies it with one of
the three torsion graphs. Hence $G$ also has nonzero integral homology.
In the second case the terminal graph is one of the $51$ stuck seeds.
A torsion seed again gives nonzero integral homology. A contractible
seed gives one of the eight graphs of
Theorem~\ref{thm:classification14}. This exhausts the possibilities.
In every torsion case the first integral homology is $\ZZ/3$ and
the reduced homology vanishes in all other degrees.

\begin{proposition}\label{prop:table-integral}
In each universe of Table~\ref{tab:card}, integral acyclicity and
I-contractibility coincide. Their common cardinality equals the
$\mathbb F_2$-acyclic count through order $12$, that count minus
three at order $13$, and that count minus $625$ in the $K_4$-free
row of order $14$.
\end{proposition}
\begin{proof}
The preceding classification partitions every acyclic non-strong
graph into a member of $\I$ or a graph with nonzero integral
homology. Together with~\eqref{eq:table-chain}, this proves the
equality of classes in the specified universes. At order $14$,
the acyclic and strong counts differ by
\[
 \RowFourteenAcyclic-\RowFourteenIS=633.
\]
Exactly eight of these graphs lie in $\I$, leaving $625$ with
nonzero integral homology. The preceding reduction also determines their homology:
$\widetilde H_1\cong\ZZ/3$, with vanishing reduced homology in
all other degrees. Separate integral calculations on the $625$
graphs confirm these groups. Thus the integral entries are
$\RowThirteenI$ at order $13$ and $\RowFourteenI$ in the last
row. All earlier entries follow from equality of the acyclic and
strong classes.
\end{proof}

The equality just proved concerns the listed finite universes; it
does not identify integral acyclicity and contractibility in general.

Table~\ref{tab:card} combines these deductions with the direct counts.
The $63$ entries shared with the published table
\cite[Table~1]{DEFH2023} agree: these are the connected,
$\mathcal A_{\mathbb F_2}$, $\I$, $\IS$, $\ISV$, $1$-Dis
and $0$-Dis columns at orders $3$ through $11$.

\begin{table}[tbp]
\centering
\caption{Numbers of isomorphism classes of connected graphs.
The two panels give twelve counts, with $n$ repeated.
The entire row $14^{\dagger}$ is restricted to $K_4$-free graphs.
$\mathcal A_R$ denotes reduced acyclicity over $R$, and $(-1)$-Dis
denotes the cones.}\label{tab:card}
% Thirteen logical columns, including n; twelve counts in two panels.
\begingroup
\footnotesize
\setlength{\tabcolsep}{2pt}
\sisetup{group-separator={,}}
\renewcommand{\arraystretch}{1.12}
\begin{tabular*}{\textwidth}{@{\extracolsep{\fill}}rrrrrrr@{}}
\toprule
$n$ & connected & $\chi=1$ & $\mathcal A_{\mathbb F_2}$ & $\mathcal A_{\ZZ}$ & $\I$ & $\IS$ \\
\midrule
3 & \num{2} & \num{2} & \num{2} & \num{2} & \num{2} & \num{2} \\
4 & \num{6} & \num{5} & \num{5} & \num{5} & \num{5} & \num{5} \\
5 & \num{21} & \num{16} & \num{16} & \num{16} & \num{16} & \num{16} \\
6 & \num{112} & \num{68} & \num{68} & \num{68} & \num{68} & \num{68} \\
7 & \num{853} & \num{404} & \num{403} & \num{403} & \num{403} & \num{403} \\
8 & \num{11117} & \num{3837} & \num{3798} & \num{3798} & \num{3798} & \num{3798} \\
9 & \num{261080} & \num{68458} & \num{66167} & \num{66167} & \num{66167} & \num{66167} \\
10 & \num{11716571} & \num{2571263} & \num{2343148} & \num{2343148} & \num{2343148} & \num{2343148} \\
11 & \num{1006700565} & \num{204131202} & \num{166978401} & \num{166978401} & \num{166978401} & \num{166978401} \\
12 & \num{164059830476} & \num{32482574568} & \num{22666537431} & \num{22666537431} & \num{22666537431} & \num{22666537431} \\
13 & \num{50335907869219} & \num{9859955384222} & \num{5619022190286} & \num{5619022190283} & \num{5619022190283} & \num{5619022190283} \\
14$^{\dagger}$ & \num{309317860269934} & \num{6112638767338} & \num{755877673752} & \num{755877673127} & \num{755877673127} & \num{755877673119} \\
\bottomrule
\end{tabular*}
\par\medskip
\begin{tabular*}{\textwidth}{@{\extracolsep{\fill}}rrrrrrr@{}}
\toprule
$n$ & $\ISV$ & $3$-Dis & $2$-Dis & $1$-Dis & $0$-Dis & $(-1)$-Dis \\
\midrule
3 & \num{2} & \num{2} & \num{2} & \num{2} & \num{2} & \num{2} \\
4 & \num{5} & \num{5} & \num{5} & \num{5} & \num{5} & \num{4} \\
5 & \num{16} & \num{16} & \num{16} & \num{16} & \num{16} & \num{11} \\
6 & \num{68} & \num{68} & \num{68} & \num{68} & \num{68} & \num{34} \\
7 & \num{403} & \num{403} & \num{403} & \num{403} & \num{403} & \num{156} \\
8 & \num{3798} & \num{3798} & \num{3798} & \num{3798} & \num{3791} & \num{1044} \\
9 & \num{66167} & \num{66167} & \num{66167} & \num{66167} & \num{65561} & \num{12346} \\
10 & \num{2343148} & \num{2343148} & \num{2343148} & \num{2343148} & \num{2258314} & \num{274668} \\
11 & \num{166978389} & \num{166978389} & \num{166978389} & \num{166978389} & \num{151084672} & \num{12005168} \\
12 & \num{22666525713} & \num{22666525713} & \num{22666525713} & \num{22666525693} & \num{18510237804} & \num{1018997864} \\
13 & \num{5619007241580} & \num{5619007241580} & \num{5619007241580} & \num{5619007205823} & \num{4015494549880} & \num{165091172592} \\
14$^{\dagger}$ & \num{755266302709} & \num{755266302709} & \num{755266302709} & \num{755266302709} & \num{81156275851} & \num{20797002} \\
\bottomrule
\end{tabular*}
\endgroup

\end{table}

\subsection{Consequences of the table}

The constructions and the complete lower-order counts now determine
the three minima relevant to the conjectures.

\begin{corollary}[Computer-assisted]\label{cor:minimality}
The minimum order of an $\mathbb F_2$-acyclic graph outside $\IS$
is $13$. The minimum order of a graph in $\I\setminus\IS$ is
$14$, and the same minimum holds for an integrally acyclic graph
outside $\IS$. The minimum order of a strong graph having an
$\IS$-admissible edge whose deletion leaves a non-strong graph
is also $14$. The minimum number of vertices of a contractible,
non-collapsible flag simplicial complex is $14$.
\end{corollary}
\begin{proof}
Theorem~\ref{thm:classification13} excludes the smaller orders,
and Theorems~\ref{thm:A} and~\ref{thm:B} provide the first two
examples. Proposition~\ref{prop:table-integral} excludes an
integrally acyclic non-strong graph through order $13$, while
$\Ga\in\I$ attains order $14$. Finally, an admissible edge
deletion from a strong graph is an I-transformation. If its result
is non-strong, that result lies in $\I\setminus\IS$ and has the
same order. The lower bound is therefore $14$; the one-edge
extension of $\Ga$ in Section~\ref{sec:examples} attains it.

If a contractible flag complex $K=\cl(G)$ has at most $13$ vertices,
then $G\in\I$ by \cite[Corollary~1.5]{DochtermannMatsushita2025}.
Theorem~\ref{thm:classification13} gives $G\in\IS$, and
\cite[Theorem~1]{EFH2022} shows that $K$ is collapsible.
Proposition~\ref{prop:noncollapsible} supplies the example
$\cl(\Ga)$ on $14$ vertices.
\end{proof}

The dismantling columns give further separations. The classes
$2$-Dis and $3$-Dis have the same counts as $\ISV$ throughout the
table. The first difference between $2$-Dis and $1$-Dis occurs at
order $12$, where it has cardinality $20$; at order $13$ the
difference is $\NumTwoDisNotOneDisThirteen$. These conclusions
follow by subtracting the counts of nested classes. The $20$ graphs
of order $12$ are catalogued in the accompanying material. The first
difference between $\IS$ and $\ISV$, already established in
\cite[Theorem~11]{DEFH2023}, occurs at order $11$ and consists of
twelve graphs.

There is also a direct interpretation of the cone column. Adjoining
a universal vertex is a bijection between isomorphism classes of
graphs on $n-1$ vertices and cones on $n$ vertices. Its inverse is
well defined because deleting any of several universal vertices
gives isomorphic graphs. This explains the agreement of the
unrestricted cone entries with the known numbers of unlabelled
graphs~\cite{OEISA000088}; the connected totals at orders $3$ through
$13$ agree with~\cite{OEISA001349}. A cone is $K_4$-free exactly when deletion
of its apex leaves a triangle-free graph. The last cone entry is
therefore $\RowFourteenCone$, the number of triangle-free graphs
on $13$ vertices~\cite{OEISA006785}.
